\chapter{Introduction}
\label{chapter:introduction}

Robots are expected to navigate in multiples and different environments, in order to achieve this we need to allow them to have situational awareness and autonomy to make proper decisions. Computer vision aims to enable machines to get an interpretation of the world that surround them. Processing images and videos robots can extract information to recognize objects geometries, textures, materials, etc. Although it is true that with visual analysis, the machine can create internal models that contain a vast knowledge of the environment, this doesn't mean that its perception is complete. Robotic perception give the ability to a machine to understand its surroundings to the degree that could safely interact with the world \cite{RobotV}, for such understanding other techniques are required as well as computer vision to get a complete sensory experience of the environment. Acoustic sensing and analysis can enhance the perception of machinery following an approach similar to that seen in nature with animals that use binaural acoustic perception, the environment can be analyzed to detect the source of a particular sound and focus the visual attention to that particular region, but, this capability is not only an extension, in some conditions the animal can entirely relies on the acoustic information to operate, this is the case of echolocation employed by bats, birds, dolphins and wales, they ear the echoes of acoustic chirps to gain spatial perception of the surrounding environment and this enable them to navigate in complex layouts \cite{echoloco}. 

 Traditional computer vision techniques that employ sensing devices such as cameras and lidars provide high-fidelity information to navigation systems, but they become unreliable in scenarios with poor light or presence dust, smoke, and fog \cite{cat1}. Another aspect to consider is the cost of the sensing device, here the acoustic perception offers an advantage due the low cost of commercial electret or MEMS microphones which are inexpensive compared to optical cameras and lidars; then become desirable its use as an alternative to traditional sensors. Recent developments in machine learning as the ones shows in \cite{bat1, cat1, echosense} demonstrated the big potential to use neural network to infer spatial and geometrical information of objects in the surrounds by just analyzing acoustic echoes in the audible spectrum. The manufacturing sector under the industry 4.0 transformation is increasing in the adoption of automation and smart monitoring technologies and so it has become an important field of application of these new methods, where the supervision of machinery during operation, especially the machines allocated in production floors using robotic arms and computer numerical control actuators intended to manipulate objects in confined spaces \cite{mordor}. In such conditions the acoustic sensing is a good approach to get an insight of the overall status of the system due to the nature of the Room Impulse Response as a transfer function that comprehends information of the environment setup. The acoustic response of a room can be use together with artificial neural networks like convolucional neural networks or recurrent neural networks to learn from this response and be capable of recognize characteristics and properties of the room echoes to recognize spaces and objects in it. Thus, acoustic perception together with neural networks can improve the perception of machinery when it operates in complex spaces.  




%\newpage

\section{Problem}

The use of acoustic signal to approximate the size and shape of objects was considerably improved by the adoption of machine learning techniques. An important part of the work in this field was develop with the scope of indoor navigation systems where the objective is to enhance the perception of robots using acoustic analysis to build low resolution deep images of an unknown environment. A different point of view is formulated and presented in this work, one in which we know the environment completely a priory, but we have an object that will change it's position inside the environment and we expect to know it's location with certain accuracy each time. Second,  contrary to the large size room in which those systems are intended to navigate, we expect to estimate the object location in short ranges in which multi-path echoes could be received at the same time creating a superposition difficult to resolve and cause malfunction in those systems \cite{bat1}. It is our understanding that there are no works related to estimating the position of an object through acoustic analysis under the circumstances described. To achieve the acoustic spatial estimation in this scenario we propose two key steps to the processing pipeline, one in the beamforming stage to increase the spatial discrimination from where the echoes are originated, and second a change in the type of neural network use to infer the location, to employ a network suitable to directly process time series instead of Convolutional Neural Networks. Our purpose is to demonstrate that is possible to locate an object in a small room (e.g. 1m side dimension) using acoustic signal processing with neural networks.    

\section{Document Structure}

This document present in \textbf{Chapter 1} a brief introduction to the topic and is stated the problem we intent to address with the development of this thesis. \textbf{Chapter 2} contains the theoretical background and state of the art in the field. \textbf{Chapter 3} details the hypothesis, objectives and scope of the work. \textbf{Chapter 4} shows the research proposal and methodology. \textbf{Chapter 5} present the schedule for the execution of the proposed work. 






